\section{研究前沿：从网页 Agent 到通用多模态执行体}

\subsection{更高层的动作抽象}
讲者提出一个值得长期投入的方向：为 Agent 设计更高层语义抽象，让模型不必在每个任务都从 click/type 原子动作开始搜索。这能显著降低搜索深度和工程复杂度。

\begin{importantbox}{高层抽象的潜在收益}
\begin{itemize}
    \item 减少低层动作噪声；
    \item 提高跨站点泛化；
    \item 为多 Agent 协同创造统一接口。
\end{itemize}
\end{importantbox}

\subsection{并行搜索与多实例执行}
讲者还提到并行化的重要性：同一任务可在多个实例上探索不同路径，再把最优轨迹回传主流程。这与大模型时代 ``test-time compute'' 增长趋势一致。

\begin{knowledgebox}{并行执行的设计要点}
\begin{enumerate}
    \item 子任务可独立回放；
    \item 共享中间评估信号；
    \item 统一冲突解决策略（例如取最高价值路径）。
\end{enumerate}
\end{knowledgebox}

\subsection{评估基准的下一步}
下一代基准需要更真实地覆盖：动态网站、个性化页面、长期会话、跨设备交互和多角色协作。否则系统会在静态 benchmark 上提升，却在真实生产环境失效。

\begin{warningbox}{Benchmark 与真实部署之间的结构性差距}
网页产品持续变化、元素动态加载、权限和会话约束复杂，任何固定 benchmark 都可能快速过时。评估体系必须支持持续更新与版本追踪。
\end{warningbox}

\subsection{本章小结}
多模态 Agent 的长期突破点在于动作抽象、并行搜索和动态评估体系。未来竞争将更多发生在系统工程层，而不仅是参数规模层。

